\documentclass[10pt]{article}
\title{Physical Mathematics}
\author{EFF}
\date{\today}

%\usepackage[a4paper]{geometry}
\usepackage[paperheight=25.5 cm ,paperwidth=14.25cm,left=1.75cm,right=1.75cm,top=1.5cm,bottom=1.5cm]{geometry}
\usepackage{amsmath, amsthm, amssymb, esdiff, esint, indentfirst, longtable, array, csquotes, fancyref, array, longtable, romannum, mathtools, booktabs, braket}
\theoremstyle{definition}
\newtheorem{definition}{Definition}[section]
\usepackage{graphicx}
\usepackage[]{subfiles}
\usepackage[]{caption}

\usepackage[backend=biber]{biblatex}
\addbibresource{./export.bib}

\graphicspath{ {./Rss/} }
\newcommand*{\df}{\frac{d}{dx}}
\newcommand*{\seconddf}{\frac{d^2}{dx^2}}
\newcommand*{\pardf}{\frac{\partial}{\partial x}}
\newcolumntype{M}[1]{>{\[}p{#1}<{\]}}
\def\rcurs{{\mbox{$\resizebox{.08in}{.08in}{\includegraphics{{../Rss/ScriptR.pdf}}}$}}}
\def\brcurs{{\mbox{$\resizebox{.08in}{.08in}{\includegraphics{{../Rss/BoldR.pdf}}}$}}}
\def\hrcurs{{\mbox{$\hat \brcurs$}}}
\newcommand{\A}{\mathbf{A}}
\newcommand{\B}{\mathbf{B}}
\newcommand{\C}{\mathbf{C}}
\newcommand{\x}{\mathbf{\hat{x}}}
\newcommand{\y}{\mathbf{\hat{y}}}
\newcommand{\z}{\mathbf{\hat{z}}}
\DeclareMathOperator{\erf}{erf}
\DeclareMathOperator{\erfc}{erfc}
\DeclareMathOperator{\erfi}{erfi}
\allowdisplaybreaks
\newcolumntype{C}{>{$}c<{$}}

\usepackage{xcolor}
\definecolor{myblue}{HTML}{2d2d3f}
\usepackage[colorlinks=true, linkcolor=myblue, urlcolor=myblue, citecolor=myblue]{hyperref}
\hypersetup{bookmarksopen=true}

\usepackage{etoolbox}
\makeatletter
\makeatother

\usepackage{titlesec}
\titleformat{\chapter}[display]
{\thispagestyle{empty}\centering \Huge \bfseries}
{}                                              
{0pt}                                             
{}                                                
\titlespacing*{\chapter}{0pt}{\fill}{\fill}

\titleformat{\section}[hang]
{\bfseries\Large\centering}
{}
{0pt}
{}
[\vspace{0.5cm}]

\titleformat{\subsection}[hang]
{\bfseries\large}
{} 
{0pt}
{}
[\noindent\rule{\textwidth}{0.4pt}]

\titleformat{\subsubsection}[runin]
{\bfseries}
{} 
{0pt}
{\indent}

\titleformat{\paragraph}[runin]
{\itshape}
{} 
{0pt}
{\indent}

\graphicspath{{./Rss/}}
\newcommand{\normfig}[1]{\includegraphics[width=0.5\textwidth]{#1}}
\newcommand{\normfigL}[1]{\includegraphics[width=0.8\textwidth]{#1}}
\newcommand{\dfig}[1]{\includegraphics[width=0.4\textwidth]{#1}}
\newcommand{\fullfig}[1]{\includegraphics[width=\textwidth]{#1}}


\begin{document}
\pagestyle{empty}
\nocite{*}
\subfile{Section/Preamble/Pre.tex}
\clearpage
\section{FUNDAMENTAL}
\thispagestyle{empty}
\subfile{Section/Fundamental/Fundamental.tex}
\clearpage
\section{Series}
\subfile{Section/Series/Series.tex}
\clearpage

\section{Complex Analysis}
\subfile{Section/ComplexNumber/ComNum.tex}
\clearpage
\subfile{Section/ComplexNumber/CNAppend.tex}
\clearpage

\section{Linear Algebra}
\subfile{Section/LinearAlgbra/LinAlg.tex}
\clearpage
\subfile{Section/LinearAlgbra/LinAlgAppl.tex}
\clearpage

\section{Partial Differentitation}
\subfile{Section/PartialDerivative/PD.tex}
\clearpage
\subfile{Section/PartialDerivative/PDAppend.tex}
\clearpage

\section{Vector Analysis}
\subfile{Section/VectorAnalysis/VA.tex}
\clearpage

\section{Fourier Series and Transform}
\subfile{Section/Fourier/FST.tex}
\clearpage

\section{Ordinary Diferential Equation}
\subfile{Section/OrdinaryDifferentialEQ/ODE.tex}
\clearpage
\subfile{Section/OrdinaryDifferentialEQ/ODEAppen.tex}
\clearpage
\subfile{Section/OrdinaryDifferentialEQ/EDEAppend2.tex}
\clearpage
\subfile{Section/OrdinaryDifferentialEQ/ODEAppend3.tex}
\clearpage

\section{Calculus of Variation}
\subfile{Section/VariationCalculus/Var.tex}
\clearpage
\subfile{Section/VariationCalculus/VarAppl.tex}
\clearpage

\section{Special Function}
\subfile{Section/SpecialFunction/Sfunc.tex}
\clearpage

\section{Probability}
\subfile{Section/Probability/Probability.tex}
\clearpage

\printbibliography
\end{document}